% Part: applied-modal-logics
% Chapter: epistemic-logic
% Section: bisimulations

\documentclass[../../../include/open-logic-section]{subfiles}

\begin{document}

\olfileid{aml}{el}{bsd}

\olsection{互模拟}

关于我们的认知语言的表达能力，可能还剩下一个问题，即模型与在其中成立的公式之间的关系。我们从框架对应性结果中已经看到，当某些公式在一个框架中有效时，它们也会使得该框架满足某些性质。但是，例如，我们的模态语言是否允许我们区分一个带有自返箭头的世界与一个无穷的世界链（其中每个世界都指向下一个）呢？也就是说，是否存在某个公式 $A$ 仅在这两个世界之一中成立？

互模拟是我们可以在关系模型之间定义的一种关系，用以说明它们实质上具有相同的结构。正如我们将要看到的，它刻画了模型之间的一种等价性，而这种等价性是能够用我们的认知语言来表达的。

\begin{defn}[互模拟]
令 $M_1 = \tuple{W_1, R_1, V_1}$ 和 $M_2 = \tuple{W_2, R_2, V_2}$ 为两个关系模型，且令 $\mathcal{R} \subseteq W_1 \times W_2$ 为一个二元关系。若对每一个 $\tuple{w_1, w_2} \in \mathcal{R}$ 满足以下条件，则称 $\mathcal{R}$ 是一个\emph{互模拟}：

\begin{enumerate}
  \item 对所有命题变元~$p$，$w_1 \in V_1(p)$ 当且仅当 $w_2 \in V_2(p)$。
  \item 对所有主体 $a \in A$ 和世界 $v_1 \in W_1$，如果 $R_{1_a} w_1 v_1$，则存在 $v_2 \in W_2$ 使得 $R_{2_a} w_2 v_2$ 且 $\tuple{v_1, v_2} \in
  \mathcal{R}$。
  \item 对所有主体 $a \in A$ 和世界 $v_2 \in W_2$，如果 $R_{2_a} w_2 v_2$，则存在 $v_1 \in W_1$ 使得 $R_{1_a} w_1 v_1$ 且 $\tuple{v_1, v_2} \in
   \mathcal{R}$。
\end{enumerate}

当 $M_1$ 和 $M_2$ 之间存在一个连接世界 $w_1$ 和 $w_2$ 的互模拟时，我们也可记作 $\tuple{M_1, w_1} \leftrightarroweq
\tuple{M_2, w_2}$，并称 $\tuple{M_1, w_1}$ 和 $\tuple{M_2, w_2}$ 是\emph{互模拟的}。
\end{defn}

互模拟关系中的不同条款各自保证了不同的方面。条款1确保互模拟的世界将满足相同的不含模态算子的公式，因为它保证了所有命题变元上的取值一致。另外两个条款（有时分别称为“前向”和“后向”条件）确保可达关系具有相同的结构。

\begin{thm}
如果 $\tuple{M_1, w_1} \leftrightarroweq \tuple{M_2, w_2}$，那么对每个公式~$!A$，$\mSat{M_1}{!A}[w_1]$ 当且仅当 $\mSat{M_2}{!A}[w_2]$。
\end{thm}

\begin{figure}
  \begin{center}
    \begin{tikzpicture}[modal]
      \node[world] (w1) {$w_1$};
      \node[world] (w2) [above left=of w1]{$w_2$};
      \node[world] (w3) [above right=of w1] {$w_3$};
      \draw[<->] (w1) to node {$a$} (w2);
      \draw[->,loop below] (w1) to node {$a$} (w1);
      \draw[<->] (w1) to [swap] node {$a$} (w3);
      \draw[->,loop above] (w2) to node {$a$} (w2) ;
      \draw[->,loop above] (w3) to node {$a$} (w3) ;

      \node[world](v1)[right of=w1, xshift=1.5in]{$v_1$};
      \node[world](v2)[above of=v1]{$v_2$};
      \draw[<->] (v1) to node {$a$} (v2);
      \draw[->,loop below] (v1) to node {$a$} (v1);
      \draw[->,loop above] (v2) to node {$a$} (v2) ;

      \draw[-,dotted] (w1) to (v1) ;
      \draw[-,dotted,bend left=45] (w2) to (v2) ;
      \draw[-,dotted,bend left=30] (w3) to (v2) ;
    \end{tikzpicture}
  \end{center}
  \caption{两个互模拟的模型。}
  \ollabel{fig:bisimilar}
\end{figure}

尽管 \olref{fig:bisimilar} 中描绘的两个模型并不完全相同，但存在一个连接世界 $w_1$ 和~$v_1$ 的互模拟。这个互模拟也将把 $w_2$ 和 $w_3$ 都连接到~$v_2$，这意味着我们的模态语言所能表达的内容无法真正区分它们。然而，如果 $w_2$ 和~$w_3$ 满足不同的命题变元，情况就会不同。

\end{document}